% Laget automatisk av elevark/lag_elevark.py fra opplegg/video-skratt-kast.html. Endringer her blir overskrevet.
\documentclass[11pt]{article}
\usepackage{elevark}
\begin{document}
\elevarkhode{Fysikk 2 · Bevegelse og krefter i to dimensjoner}{Videoanalyse av skrått kast}{https://fysikklab.no/opplegg/video-skratt-kast.html}{Opplegget på nett}

Et skrått kast er to bevegelser på én gang: jevn fart bortover og fritt fall opp og ned. Film et kast, og la Tracker gi dere $x$ og $y$ for hvert bilde. Stemmer det at $v_{x}$ er konstant og at $a_{y}$ = −g?

\section*{Hensikt}

\begin{itemize}
\item Analysere et skrått kast som to uavhengige bevegelser i x- og y-retning.
\item Vise at farten i x-retning er konstant, og at akselerasjonen i y-retning er $g$ nedover.
\item Bestemme startfarten og kastvinkelen, og forutsi hvor ballen lander.
\end{itemize}

\section*{Utstyr}

\begin{itemize}
\item En liten, tung ball, for eksempel en tennisball eller sprettball.
\item Mobil eller kamera på stativ, helst med 60 bilder per sekund.
\item En meterstokk som målestokk, og helst en vegg med ensfarget bakgrunn.
\item Tracker (gratis fra physlets.org/tracker).
\end{itemize}

\section*{Framgangsmåte}

\begin{enumerate}
\item Sett kameraet på stativ 3–4~m fra veggen, slik at hele banen kommer med. Kameraet skal stå vinkelrett på veggen, med linsa omtrent midt i banen.
\item Legg meterstokken langs veggen i samme avstand fra kameraet som kastet. Den er målestokken i Tracker.
\item Kast ballen skrått oppover langs veggen, ikke mot eller fra kameraet. Et lett kast på 1,5–2~m er nok. Film flere kast og velg det beste.
\item I Tracker: Sett kalibreringsstaven på meterstokken og plasser origo der ballen slippes, med $y$ oppover. Klikk på midten av ballen i hvert bilde fra den slipper hånden til like før den lander.
\item Kopier kolonnene $t$, $x$ og $y$, og lim dem inn i verktøyet på nettsiden. Analyser $x$ og $y$ hver for seg.
\end{enumerate}

Uten egen video kan dere bruke eksempeldataene, eller klikke dere gjennom en virtuell video av et kast i Videoanalyse uten video\footnote{\url{fysikklab.no/sim/videoanalyse.html}}. Velg «Ball som kastes skrått».

\section*{Oppgaver}

\begin{oppgaveliste}
\oppgave \fase{gjett}{Gjett}Skisser hvordan $x(t)$ og $y(t)$ ser ut for et skrått kast. Hvordan ser $v_{x}(t)$ og $v_{y}(t)$ ut?
\svar
\oppgave \fase{test}{Mål}Film et kast og analyser det i Tracker, eller bruk eksempeldataene. Er $x(t)$ en rett linje? Hva er $v_{x}$?
\svar
\oppgave \fase{test}{Mål}Tilpass en andregradsfunksjon til $y(t)$. Hva blir $a_{y}$, og hvor godt stemmer det med $g = 9{,}81\,\mathrm{m/s^{2}}$? Kryss av for «Vis fart» og se på $v_{y}(t)$.
\svar
\oppgave \fase{test}{Regn}Bruk $v_{x}$ og $v_{0}y$ til å finne startfarten og kastvinkelen. Når er ballen på toppen, og hvor høyt over utgangspunktet kommer den? Sjekk med dataene.
\svar
\oppgave \fase{forklar}{Forklar}Hvorfor er $v_{x}$ konstant gjennom hele kastet? Hvordan ville $x(t)$ sett ut for en lett ball, for eksempel en ballong, med mye luftmotstand?
\svar
\oppgave \fase{forklar}{Vurder}Ballen er nærmere kameraet enn meterstokken hvis dere ikke kaster langs veggen. Blir målt fart og akselerasjon da for store eller for små? Hvilke andre feilkilder er det?
\svar
\oppgave \fase{utfordring}{Utfordring}Bruk startfarten, kastvinkelen og høyden der ballen slipper, til å regne ut hvor langt bortover ballen kommer før den treffer gulvet. Sammenlign med videoen. I eksemplet slipper ballen 0,85~m over gulvet.
\svar
\end{oppgaveliste}

\end{document}
